\section{Symbols, conventions, and terminology}
\label{app:notation}

We write $\sym{A}$ for a reference operator and
$\sym{Acurrent}=\sym{A}+\sym{H}$ for the perturbed operator. The
sine-theta theorem uses coordinate maps for trial and exact subspaces.
The transport results name the subspaces $\sym{Usource},\sym{Vtarget}$
and their projections $\sym{Psource},\sym{Qtarget}$ directly.
The sine-theta projection is $\sym{Pexact}$, the residual is $\sym{R}$,
and direct rotations are denoted by $\sym{Rrotation}$ throughout.

In the colored version, blue identifies exact or target data, amber trial
or source data, cyan the exact complementary block, magenta angles, green
residuals, and violet spectral separation. The residual $\sym{R}$ and
direct rotation $\sym{Rrotation}$ have distinct roles and colors. The
notation and definitions also apply in monochrome.

\begingroup
\small
\setlength{\tabcolsep}{4pt}
\renewcommand{\arraystretch}{1.15}
\begin{longtable}{@{}>{\raggedright\arraybackslash}p{0.29\textwidth}>{\raggedright\arraybackslash}p{0.65\textwidth}@{}}
\toprule
\textbf{Symbol} & \textbf{Definition or interpretation}\\
\midrule
\endfirsthead
\multicolumn{2}{@{}l}{\textit{Symbols and conventions (continued)}}\\[2pt]
\toprule
\textbf{Symbol} & \textbf{Definition or interpretation}\\
\midrule
\endhead
\bottomrule
\endfoot
$\sym{A},\sym{H},\sym{Acurrent}$ & Reference operator, additive perturbation, and perturbed operator $\sym{Acurrent}=\sym{A}+\sym{H}$.\\
$\sym{AmbientHilbertSpace},\ \sym{TrialCoordinateSpace},\ \sym{ExactCoordinateSpace},\ \sym{ComplementaryCoordinateSpace}$ & Ambient, trial-coordinate, exact-coordinate, and complementary-coordinate Hilbert spaces.\\
$\sym{E0}:\sym{TrialCoordinateSpace}\to\sym{AmbientHilbertSpace}$ & \dkSymbolDescription{E0}\\
$\sym{TrialOperator}$ & Self-adjoint trial operator. A Rayleigh--Ritz choice is $\sym{TrialOperator}=\like{E0}{\symglyph{E0}^*}\sym{Acurrent}\sym{E0}$.\\
$\sym{ReferenceBlock0},\sym{ReferenceBlock1}$ & Reference blocks on $\DKRange(\sym{Psource})$ and $\DKKernel(\sym{Psource})$ in Section~\ref{sec:foundations}.\\
$\sym{F0}:\sym{ExactCoordinateSpace}\to\sym{AmbientHilbertSpace}$ & \dkSymbolDescription{F0}\\
$\sym{F1}:\sym{ComplementaryCoordinateSpace}\to\sym{AmbientHilbertSpace}$ & \dkSymbolDescription{F1}\\
$\sym{Lambda1}$ & Complementary self-adjoint block, with $\sym{Acurrent}\sym{F1}=\sym{F1}\sym{Lambda1}$ on its domain.\\
$\sym{SpectralMeasure},\sym{SpectralSet}$ & Projection-valued measure and Borel spectral set; $\sym{SpectralMeasure}(\sym{SpectralSet})$ is the corresponding spectral projection. These symbols differ from the ambient space $\sym{AmbientHilbertSpace}$ and complementary block $\sym{S}$.\\
$\sym{Pexact}=\sym{F0}\like{F0}{\symglyph{F0}^*}$ & Projection onto the desired exact subspace in the sine-theta theorem.\\
$\sym{Theta0},\sym{sinTheta0}$ & Trial-coordinate angle operator and its positive sine, with $\sym{sinTheta0}=(\sym{Identity}-\like{E0}{\symglyph{E0}^*}\sym{Pexact}\sym{E0})^{1/2}$.\\
$\sym{Theta}=\sym{DiagonalOperator}(\sym{Theta0},\sym{Theta1})$ & Full angle operator; $\sym{Theta1}$ compares the orthogonal complements. Its operator norm is the largest spectral angle.\\
$\sym{CrossedSourceDefect},\sym{CrossedTargetDefect}$ & Source directions perpendicular to the target, and target directions perpendicular to the source. Section~\ref{sec:foundations} uses the condition $\sym{DimensionOperator}\sym{CrossedSourceDefect}=\sym{DimensionOperator}\sym{CrossedTargetDefect}$.\\
$\{\sym{ZeroVector}\}$ & Zero subspace: the set containing only the zero vector.\\
$\sym{S}=(\sym{Identity}-\sym{Pexact})\sym{E0}$ & Complementary block, with modulus $\dkModulus{\sym{S}}=\sym{sinTheta0}$.\\
$\sym{X}=\like{F1}{\symglyph{F1}^*}\sym{E0}$ & Coordinates of the trial component in the unwanted space; $\sym{S}=\sym{F1}\sym{X}$.\\
$\sym{SylvesterLeft},\sym{SylvesterRight},\sym{SylvesterRHS}$ & Generic left and right Sylvester coefficients and right side: $\sym{SylvesterLeft}\sym{X}-\sym{X}\sym{SylvesterRight}=\sym{SylvesterRHS}$. The historical equation writes the right coefficient as $\sym{HistoricalSylvesterRight}$.\\
$\sym{R}=\sym{Acurrent}\sym{E0}-\sym{E0}\sym{TrialOperator}$ & Bounded residual agreeing with this difference on $\DKDom(\sym{TrialOperator})$.\\
$\sym{delta}>0$ & \dkSymbolDescription{delta}\\
$[\sym{betaEndpoint},\sym{alphaEndpoint}]$ & Interval in the interval/exterior gap hypothesis.\\
$\DKRange(T),\DKKernel(T),\DKDom(T)$ & Range, kernel, and domain of an operator.\\
$\sym{Usource},\sym{Vtarget}$ & Source and target subspaces in the displacement problem.\\
$\sym{Psource},\sym{Qtarget}$ & Orthogonal projections onto $\sym{Usource},\sym{Vtarget}$ in that problem.\\
$\sym{Rrotation},\sym{Wcompetitor}$ & Direct rotation and an arbitrary orthogonal or unitary transporter carrying $\sym{Usource}$ onto $\sym{Vtarget}$.\\
$\sym{PolarU}$ & Partial isometry in a polar decomposition, as in $\sym{S}=\sym{PolarU}|\sym{S}|$.\\
$\sym{IntertwinerT}$ & The intertwiner $\sym{Qtarget}\sym{Psource}+(\sym{Identity}-\sym{Qtarget})(\sym{Identity}-\sym{Psource})$, whose unitary polar factor is the acute direct rotation.\\
$\sym{theta},\like{theta}{\symglyph{theta}_j}$ & Individual principal angles, in $[0,\sym{PiConstant}/2]$.\\
$\sym{Xframe},\sym{Yframe}$ & Orthonormal coordinate matrices for the source and target in the real counterexample.\\
$\sym{Jturn},\sym{Cturn},\sym{Hmix}$ & Planar rotation through $\sym{PiConstant}/2$, rotation through $2\sym{theta}$, and orthogonal mixing matrix used in that counterexample.\\
$\sym{Jreflect};\ \sym{AngleGenerator}$ & Source reflection $2\sym{Psource}-\sym{Identity}$, with square $\sym{Identity}$; skew angle generator of a direct rotation. The $\sym{PiConstant}/2$-rotation matrix $\sym{Jturn}$ has square $-\sym{IdentityTwo}$.\\
$\sym{CounterPrincipalOne},\sym{CounterPrincipalTwo};\ \sym{CounterMovingPlane},\sym{CounterFixedPlane}$ & Principal planes of the direct rotation and moving/fixed planes of the competitor in \cref{fig:prop44-displacement-planes}.\\
$\sym{CounterChordA},\sym{CounterChordB}$ & Displacement singular values $\sqrt{2-\sqrt2}$ and $\sqrt2$ in that counterexample.\\
$\sym{Kdisplacement}=\sym{Identity}-\sym{Wcompetitor}$ & Full displacement operator in the failed Ky Fan comparison.\\
$\sym{OmegaOne},\sym{OmegaTwo}$ & Principal-plane projections used to test that comparison.\\
$\dkModulus{T}=(T^*T)^{1/2}$ & Positive modulus of an operator.\\
$\sym{SingularValue}(T)$ & Ordered singular values; approximation numbers provide the corresponding comparisons for bounded operators in Hilbert space.\\
$\sym{NormN}(T)$ & \dkSymbolDescription{NormN}\\
$\dkOperatorNorm{T}$ & Operator norm; in finite dimension, the largest singular value.\\
$\dkFrobeniusNorm{T}$ & Frobenius norm, $\bigl(\sum_j \sym{SingularValue}(T)^2\bigr)^{1/2}$.\\
$\dkTraceNorm{T}$ & Trace norm, $\sum_j \sym{SingularValue}(T)$.\\
$\dkKyFanNorm{T}{j}$ & Ky Fan norm, the sum of the $j$ largest singular values.\\
$\sym{NormQ}\text{-norm}$ & A norm satisfying $\sym{NormN}(T)^2=\sym{NormM}(T^*T)$ for a unitarily invariant norm $\sym{NormM}$.\\
$\sym{Matching},\like{SingletonCost}{\symglyph{SingletonCost}_i},\like{PairCost}{\symglyph{PairCost}_{ij}}$ & Matching of positive principal-angle labels, with singleton and paired costs in \cref{eq:nuclear-matching-costs}.\\
$\sym{MatchingEdgeWeight};\ \like{TraceCost}{\symglyph{TraceCost}(M)}$ & Edge weight $\like{SingletonCost}{\symglyph{SingletonCost}_i}+\like{SingletonCost}{\symglyph{SingletonCost}_j}-\like{PairCost}{\symglyph{PairCost}_{ij}}$; total trace-norm displacement cost of the matching.\\
\end{longtable}
\endgroup

\subsection{Terminology}
\label{app:terminology}

The following terms describe the roles of the mathematical objects used in
\cref{sec:sine-theta,sec:example-transport}.

\begin{description}
\item[Ambient space and operator.] The whole Hilbert space
$\sym{AmbientHilbertSpace}$ in which the subspaces lie, and the operator acting
on that space. In the sine-theta statement the ambient operator is
$\sym{Acurrent}$.

\item[Reference and perturbed operators.] The reference operator is
$\sym{A}$; an additive perturbation $\sym{H}$ gives the perturbed operator
$\sym{Acurrent}=\sym{A}+\sym{H}$. The choice of which operator supplies the
desired spectral subspace depends on the approximation problem.

\item[Exact subspace.] The desired subspace specified by the spectral
decomposition of the operator of interest. In the sine-theta statement,
it is $\DKRange(\sym{F0})$ for $\sym{Acurrent}$, with orthogonal complement
$\DKRange(\sym{F1})$. In the perturbation application in the introduction,
the desired subspace instead belongs to the reference operator $\sym{A}$.
An exact formula, as in the transport theorem, gives an equality for the
quantity being determined.

\item[Trial subspace and trial operator.] The trial subspace is a candidate
approximation to the desired subspace, represented by $\DKRange(\sym{E0})$.
The self-adjoint trial operator $\sym{TrialOperator}$ represents an
approximation to the action of the ambient operator in trial coordinates.

\item[Residual.] The difference
$\sym{Acurrent}\sym{E0}-\sym{E0}\sym{TrialOperator}$ between the ambient
action and the embedded trial action. For unbounded operators this identity
holds on the appropriate domain, with a bounded residual $\sym{R}$.

\item[Spectral separation.] A lower bound $\sym{delta}>0$ on the distance
between the relevant spectra. In the sine-theta theorem these are the
spectra of the trial operator and the complementary block $\sym{Lambda1}$;
the interval and ordered separation hypotheses are given in
\cref{sec:gap-geometry}.

\item[Transporter and competitor.] A transporter is an ambient orthogonal
or unitary map carrying the source subspace onto the target subspace.
A competitor is another map satisfying the same constraint.

\item[Full displacement.] The operator $\sym{Identity}-\sym{Wcompetitor}$
on the whole ambient space for a transporter $\sym{Wcompetitor}$. Its
restriction to source directions is represented by
$(\sym{Identity}-\sym{Wcompetitor})\sym{Psource}$.
\end{description}

\section{Reading the Lean signatures}
\label{app:lean-guide}
\label{app:reading-lean}

A Lean theorem declaration has the form
\[
  \texttt{theorem name (parameters) (hypotheses) : conclusion := proof}.
\]
Parameters introduce the mathematical objects under discussion.  A hypothesis
such as \texttt{(hA : IsSelfAdjoint A)} supplies the corresponding fact;
bundled predicates such as \texttt{IsTrialResidual} and
\texttt{IsExactSpectralDecomposition} collect related assumptions.  Braces mark
arguments normally inferred by Lean, while square brackets request typeclass
instances supplying ambient algebraic, analytic, or topological structure.  The proposition after the final colon is the theorem's conclusion
\citep{demoura2021lean4}.

In \Cref{lst:dk-sintheta-headline}, \texttt{RCLike} is the Mathlib abstraction
covering $\sym{RealField}$ and $\sym{ComplexField}$.  A type
\texttt{ContinuousLinearMap K F E} is a continuous linear map.  A type
\texttt{LinearPMap K E E} is a linear partial map whose
domain is a submodule; this is the representation used for possibly unbounded
operators.  \texttt{N.Mem T} states that the ideal gauge associated with $\sym{NormN}$ is
finite on $T$, and \texttt{N.gaugeReal T} is the corresponding real value.  The
suffix \texttt{Real} refers to the codomain of the gauge and not to the scalar
field.

\subsection{Paper notation and Lean binder names}
\label{app:paper-lean-notation}
The following correspondence is for the particular sine-theta theorem
in \Cref{lst:dk-sintheta-headline}, not a global naming rule for the Lean
library. The Lean source is
\texttt{DavisKahan/Sources/DavisKahan1970/SineTheta/Presentation.lean},
the declaration \texttt{sinTheta\_unbounded\_formGap\_whereDefinedUIN\_rclike}.
Lean's exact source uses Unicode subscripts and Greek letters (for example,
its trial-operator binder is typeset here as $\mathtt{A_0}$). The listing
in the main text uses ASCII transliterations and readable type names.
In particular, the Lean binder $\mathtt{A_0}$ denotes the general
$\sym{TrialOperator}$, while $\sym{ReferenceBlock0}$ is the separate
reference spectral block and does not occur as a parameter of this theorem.
The variable $\mathtt{H}$ is a coordinate-space carrier, not the
perturbation $\sym{H}$. Lean's scalar-field type variable is rendered as \texttt{K} in the
schematic listing (the source uses a blackboard-bold k).

\begingroup
\small
\setlength{\tabcolsep}{4pt}
\renewcommand{\arraystretch}{1.11}
\begin{longtable}{@{}>{\centering\arraybackslash}p{0.22\linewidth}>{\centering\arraybackslash}p{0.16\linewidth}>{\raggedright\arraybackslash}p{0.55\linewidth}@{}}
\caption{Paper symbols and the corresponding Lean binders in the sine-theta declaration.}
\label{tab:paper-lean-notation}\\
\toprule
This paper & Lean binder & Mathematical role \\
\midrule
\endfirsthead
\toprule
This paper & Lean binder & Mathematical role \\
\midrule
\endhead
\bottomrule
\endfoot
\dkLeanSineThetaRows
\end{longtable}
\endgroup

\section{Formalized Davis--Kahan results}
\label{app:dk-inventory}

Table~\ref{tab:dk-full-inventory} lists the \DKResultCount{} source targets:
the four headline theorems \citep[Section~2]{davis1970rotation} and every named theorem,
proposition, lemma, and corollary established in the paper. The table keeps
compact IDs for reference. The artifact maps each ID to the corresponding
Lean declarations.
\ifDKAnonymous\else
The \href{https://github.com/AIQ-Kitware/aiq-dkps-formalization/blob/main/dev/davis-kahan-1970-formalization-result-inventory.json}{result-to-Lean mapping}
contains that correspondence.
\fi

\begin{longtable}{>{\raggedright\arraybackslash}p{0.24\linewidth}>{\raggedright\arraybackslash}p{0.68\linewidth}}
\caption{The \DKResultCount{} Davis--Kahan results selected for formalization.}
\label{tab:dk-full-inventory}\\
\toprule
ID & Result \\
\midrule
\endfirsthead
\toprule
ID & Result \\
\midrule
\endhead
\path{S2-sin-theta} & Single-angle sine theorem \\
\path{S2-tan-theta} & Single-angle tangent theorem \\
\path{S2-sin-two-theta} & Double-angle sine theorem \\
\path{S2-tan-two-theta} & Double-angle tangent theorem \\
\path{DK-3.1-prop} & Acute direct rotation existence and uniqueness \\
\path{DK-3.2-prop} & Nonacute existence criterion \\
\path{DK-3.3-prop} & Principal square-root characterization \\
\path{DK-3.4-prop} & Square as a direct rotation \\
\path{DK-3.1-thm} & Classification of pairs of subspaces \\
\path{DK-3.1-cor} & Compact classification by angle eigenvalues \\
\path{DK-3.5-prop} & Angle commutation and eigenspace geometry \\
\path{DK-3.2-cor} & Reversal symmetry \\
\path{DK-4.1-prop} & Pointwise and singular-value extremality of the direct rotation \\
\path{DK-4.1-cor} & Unitarily invariant norm minimality of restricted displacement \\
\path{DK-4.2-prop} & Basis-angle square-sum extremality \\
\path{DK-4.3-prop} & Squared-displacement unitarily invariant norm minimality \\
\path{DK-4.4-prop} & Full-displacement optimality claim \\
\path{DK-5.1-thm} & Banach-space Sylvester lower bound \\
\path{DK-5.2-thm} & Semibounded self-adjoint Sylvester theorem \\
\path{DK-5.1-lem} & Strong-cutoff convergence of singular values \\
\path{DK-6.1-lem} & Direct-sum norm comparison and converse \\
\path{DK-6.2-lem} & Reflection-pinch contraction \\
\path{DK-6.1-prop} & Sine proof, ambient limitation, and symmetric sine theorem \\
\path{DK-6.1-thm} & Generalized sine theorem \\
\path{DK-6.2-thm} & Pairwise-gap square-norm sine theorem \\
\path{DK-6.3-thm} & Generalized tangent theorem and supporting machinery \\
\path{DK-6.3-lem} & Finite-rank near-maximizer off-diagonal block estimate \\
\path{DK-8.1-thm} & Branch selection and spectral repulsion \\
\path{DK-8.2-thm} & Smallness selects the acute branch \\
\bottomrule
\end{longtable}

A single source result may use several Lean declarations. The IDs identify
source statements independently of the organization of the formal library.
The printed claim \citep[Proposition~4.4]{davis1970rotation}
is refuted by the counterexample in
Section~\ref{sec:prop44}; its Q-norm consequence is stated separately there.

% Detailed table following the compact summary in app:dk-inventory.
% Summaries checked against the registered Lean statements and their supporting
% declarations; see TABLE4_LEAN_AUDIT.md for the row-by-row comparison.
\subsection{Hypotheses and conclusions}
\label{app:dk-result-summaries}

Table~\ref{tab:dk-result-summaries} summarizes the hypotheses and main
conclusions of the \DKResultCount{} results listed in
Table~\ref{tab:dk-full-inventory}, using the same IDs
\citep[Sections~2--8]{davis1970rotation}.
Hilbert spaces are separable and real or complex unless specified otherwise.
$\sym{NormN}$ denotes an arbitrary unitarily invariant norm on its domain of
finiteness; each displayed norm is evaluated on an operator in that domain.
$\dkOperatorNorm{\cdot}$ and $\dkHilbertSchmidtNorm{\cdot}$ denote the operator
and Hilbert--Schmidt norms. Self-adjoint operators may be unbounded:
operator identities hold on their domains, and residuals or right-hand sides
in norm estimates have bounded extensions.
For the ambient tangent bound, the tangent is defined as a bounded operator,
equivalently $\dkOperatorNorm{\sym{Psource}-\sym{Qtarget}}<1$.

For the geometric results, $\sym{Psource},\sym{Qtarget}$ are orthogonal
projections, $\sym{SourceDirectRotation}$ is a direct rotation, and
$\sym{Wcompetitor}$ is any orthogonal or unitary transporter satisfying
$\sym{Wcompetitor}\sym{Psource}\sym{Wcompetitor}^*=\sym{Qtarget}$.
Write
\[
 \sym{CrossedSourceDefect}=\DKRange(\sym{Psource})\cap\DKKernel(\sym{Qtarget}),\qquad
 \sym{CrossedTargetDefect}=\DKKernel(\sym{Psource})\cap\DKRange(\sym{Qtarget}).
\]
The pair is acute when $\sym{CrossedSourceDefect}=\sym{CrossedTargetDefect}=\{\sym{ZeroVector}\}$.
For geometric and displacement rows from \path{DK-3.3-prop} onward,
the ranges and kernels of
$\sym{Psource},\sym{Qtarget}$ have matching dimensions and
$\sym{DimensionOperator} \sym{CrossedSourceDefect}=\sym{DimensionOperator} \sym{CrossedTargetDefect}$; the displacement results also assume
the directed sine block
$(\sym{Identity}-\sym{Qtarget})|_{\DKRange(\sym{Psource})}$ compact.
$\sym{Theta0},\sym{Theta1}$ are the source and complementary angle operators,
$\sym{Theta}$ their ambient direct sum, and $\like{theta}{\symglyph{theta}_j}$ the ordered principal
angles. In the source decomposition the direct rotation has blocks
$\bigl(\begin{smallmatrix}C_0&-S_0^*\\S_0&C_1\end{smallmatrix}\bigr)$
with $C_j\ge0$; $\sym{AngleGenerator}$ is its skew angle generator, so
$\sym{SourceDirectRotation}=e^{\sym{AngleGenerator}\sym{Theta}}$.
The symbol $\sym{SourceDirectRotation}$ includes nonacute direct rotations;
for an acute pair it is the unique rotation denoted by the same letter
in the main text.
For real unitaries, spectral half-plane conditions refer to their
complexifications.

For perturbation results, $\sym{TrialOperator}$ is a self-adjoint trial operator.
In a reducing decomposition of the reference operator $\sym{A}$, the
source and complementary reference blocks are $\sym{ReferenceBlock0}$
and $\sym{ReferenceBlock1}$. The trial operator may be chosen to equal the
source reference block when the trial coordinates identify that subspace.
The operators $\sym{Lambda0},\sym{Lambda1}$ are the exact blocks of
$\sym{Acurrent}=\sym{A}+\sym{H}$.
The maps $\sym{E0},\sym{F0}$ give source and target coordinates,
$\sym{H0},\sym{H1}$ are the diagonal perturbation blocks, and
$\sym{R}=\sym{Acurrent}\sym{E0}-\sym{E0}\sym{TrialOperator}$.
Coordinate maps are isometries except in the two lower-frame rows.
The exact frames have orthogonal ranges whose sum is the ambient space;
reducing projections preserve operator domains.
We abbreviate the two spectral hypotheses, for $\sym{delta}>0$, by
\begin{align*}
 \sym{IntervalSeparation}_{\sym{delta}}(B,C):\quad&
 \sym{SpectrumSymbol}(B)\subset[\sym{betaEndpoint},\sym{alphaEndpoint}],\quad
 \sym{SpectrumSymbol}(C)\cap(\sym{betaEndpoint}-\sym{delta},\sym{alphaEndpoint}+\sym{delta})=\varnothing,
 \quad\text{or the roles interchanged};\\
 \sym{OrderedSeparation}_{\sym{delta}}(B,C):\quad&
 B\le\sym{alphaEndpoint} \sym{Identity},\quad C\ge(\sym{alphaEndpoint}+\sym{delta})\sym{Identity}.
\end{align*}
The endpoints are real, satisfy $\sym{betaEndpoint}\le\sym{alphaEndpoint}$, and are chosen for each row.
The ordered condition $\sym{OrderedSeparation}_{\sym{delta}}(B,C)$ places the spectra in
the half-infinite intervals $(-\infty,\sym{alphaEndpoint}]$ and
$[\sym{alphaEndpoint}+\sym{delta},\infty)$; it allows both operators to be unbounded.
To include both orientations, set
\[
 \sym{SineGap}(B,C)\quad\Longleftrightarrow\quad
 \sym{IntervalSeparation}_{\sym{delta}}(B,C)\quad\text{or}\quad
 \sym{OrderedSeparation}_{\sym{delta}}(B,C)\quad\text{or}\quad
 \sym{OrderedSeparation}_{\sym{delta}}(C,B).
\]
This is the separation condition of Section~\ref{sec:gap-geometry}.

\begingroup
\small
\setlength{\tabcolsep}{4pt}
\renewcommand{\arraystretch}{1.15}
\begin{longtable}{@{}>{\raggedright\arraybackslash}p{0.20\linewidth}>{\raggedright\arraybackslash}p{0.34\linewidth}>{\raggedright\arraybackslash}p{\dimexpr0.46\linewidth-4\tabcolsep\relax}@{}}
\caption{Hypotheses and main conclusions of the \DKResultCount{} Davis--Kahan results.}
\label{tab:dk-result-summaries}\\
\toprule
\textbf{ID} & \textbf{Hypotheses} & \textbf{Conclusion}\\
\midrule
\endfirsthead
\multicolumn{3}{@{}l}{\textit{Hypotheses and conclusions (continued)}}\\[2pt]
\toprule
\textbf{ID} & \textbf{Hypotheses} & \textbf{Conclusion}\\
\midrule
\endhead
\midrule
\multicolumn{3}{r@{}}{\textit{Continued on the next page}}\\
\endfoot
\bottomrule
\endlastfoot
\multicolumn{3}{@{}l}{\textbf{Headline perturbation bounds}}\\*\addlinespace[3pt]
\path{S2-sin-theta} &
 $\sym{SineGap}(\sym{TrialOperator},\sym{Lambda1})$. &
 $\sym{delta}\sym{NormN}(\sin\sym{Theta0})\le\sym{NormN}(\sym{R})$.\\\addlinespace
\path{S2-tan-theta} &
 $\sym{OrderedSeparation}_{\sym{delta}}(\sym{TrialOperator},\sym{Lambda1})$;
 Rayleigh--Ritz condition $\sym{H0}=0$. &
 $\sym{delta}\sym{NormN}(\tan\sym{Theta0})\le\sym{NormN}(\sym{R})$;
 $\sym{delta}\sym{NormN}(\tan\sym{Theta})\le\sym{NormN}(\sym{H})$.\\\addlinespace
\path{S2-sin-two-theta} &
 $\sym{A},\sym{Acurrent}$ have a common domain;
 $\sym{SpectrumSymbol}(\sym{Lambda0})\subset[\sym{betaEndpoint},\sym{alphaEndpoint}]$;
 $\sym{SpectrumSymbol}(\sym{Lambda1})$ outside
 $(\sym{betaEndpoint}-\sym{delta},\sym{alphaEndpoint}+\sym{delta})$. &
 $\sym{delta}\sym{NormN}(\sin(2\sym{Theta0}))\le2\sym{NormN}(\sym{R})$;
 $\sym{delta}\sym{NormN}(\sin(2\sym{Theta}))\le2\sym{NormN}(\sym{H})$.\\\addlinespace
\path{S2-tan-two-theta} &
 $\sym{OrderedSeparation}_{\sym{delta}}(\sym{TrialOperator},\sym{A1})$;
 bounded $\sym{H}$, $\sym{H0}=\sym{H1}=0$. &
 $\sym{delta}\sym{NormN}(\tan(2\sym{Theta0}))\le2\sym{NormN}(\sym{R})$;
 $\sym{delta}\sym{NormN}(\tan(2\sym{Theta}))\le2\sym{NormN}(\sym{H})$.\\\addlinespace[5pt]
\multicolumn{3}{@{}l}{\textbf{Two subspaces and direct rotations}}\\*\addlinespace[3pt]
\path{DK-3.1-prop} &
 $\sym{CrossedSourceDefect}=\sym{CrossedTargetDefect}=\{\sym{ZeroVector}\}$. &
 A unique direct rotation exists. Positive diagonal blocks characterize it
 among transporters.\\\addlinespace
\path{DK-3.2-prop} &
 General pair of projections. &
 A direct rotation exists iff $\sym{DimensionOperator} \sym{CrossedSourceDefect}=\sym{DimensionOperator} \sym{CrossedTargetDefect}$.
 On $\sym{CrossedSourceDefect}\oplus \sym{CrossedTargetDefect}$, $\sym{SourceDirectRotation}^2=-\sym{Identity}$;
 uniqueness can fail.\\\addlinespace
\path{DK-3.3-prop} &
 A direct rotation $\sym{SourceDirectRotation}$; conversely a unitary
 square root with spectrum in the closed right half-plane that maps
 $\sym{CrossedSourceDefect}$ onto $\sym{CrossedTargetDefect}$. &
 $\sym{SourceDirectRotation}^2=(2\sym{Qtarget}-\sym{Identity})(2\sym{Psource}-\sym{Identity})$,
 and its spectrum lies in the closed right half-plane.
 Conversely, such a square root is a direct rotation.\\\addlinespace
\path{DK-3.4-prop} &
 $C_0^2\ge\tfrac12 \sym{Identity}$;
 $Q_-=(2\sym{Psource}-\sym{Identity})\sym{Qtarget}(2\sym{Psource}-\sym{Identity})$. &
 $\sym{SourceDirectRotation}^2$ is a direct rotation from
 $\DKRange(Q_-)$ to $\DKRange(\sym{Qtarget})$.\\\addlinespace
\path{DK-3.1-thm} &
 Matching dimensions as above; angle operators
 $0\le\sym{Theta0},\sym{Theta1}\le\sym{PiConstant}/2$. &
 The four intersection dimensions and the generic angle's spectral
 multiplicity completely classify pairs up to isometric equivalence.
 Any angle data with matching nonzero multiplicities
 and the prescribed total dimension are realizable.\\\addlinespace
\path{DK-3.1-cor} &
 Classification hypotheses;
 $\sym{Psource}(\sym{Identity}-\sym{Qtarget})\sym{Psource}$ compact. &
 Classification reduces to angle eigenvalues with multiplicity:
 $\like{theta}{\symglyph{theta}_j}\downarrow0$. The two angle operators may have different
 zero multiplicities.\\\addlinespace
\path{DK-3.5-prop} &
 A direct rotation and its angle operator $\sym{Theta}$; the maximal
 eigenspace characterization assumes the acute case. &
 $\sym{Theta}$ commutes with $\sym{Psource},\sym{Qtarget},\sym{AngleGenerator},
 \sym{SourceDirectRotation}$. If $x\ne0$ and
 $\sym{Theta}x=\sym{theta} x$, then
 $\angle(x,\sym{SourceDirectRotation}x)=\sym{theta}$.
 In the acute case these eigenspaces are the unique maximal reducing
 subspaces of the prescribed angle.\\\addlinespace
\path{DK-3.2-cor} &
 Interchange $\sym{Psource}$ and $\sym{Qtarget}$. &
 $\sym{Theta}$ is unchanged, $\sym{AngleGenerator}$ changes to $-\sym{AngleGenerator}$, and the corresponding
 direct rotation becomes $\sym{SourceDirectRotation}^*$.\\\addlinespace[5pt]
\multicolumn{3}{@{}l}{\textbf{Displacement extremality}}\\*\addlinespace[3pt]
\path{DK-4.1-prop} &
 Compact angle data; arbitrary transporter $\sym{Wcompetitor}$. &
 For the positive principal angles, there are orthonormal
 $v_j\in\DKRange(\sym{Psource})$ with
 $\angle(v_j,\sym{Wcompetitor}v_j)\ge\like{theta}{\symglyph{theta}_j}$.
 Each singular value of restricted displacement is minimized by
 $\sym{SourceDirectRotation}$, with value $2\sin(\like{theta}{\symglyph{theta}_j}/2)$.\\\addlinespace
\path{DK-4.1-cor} &
 Same displacement setting. &
 $\sym{NormN}((\sym{Identity}-\sym{SourceDirectRotation})\sym{Psource})
 \le\sym{NormN}((\sym{Identity}-\sym{Wcompetitor})\sym{Psource})$.\\\addlinespace
\path{DK-4.2-prop} &
 Any orthonormal basis $(v_j)$ of $\DKRange(\sym{Psource})$;
 arbitrary transporter $\sym{Wcompetitor}$. &
 $\sum_j\sin^2\angle(v_j,\sym{Wcompetitor}v_j)
 \ge\sum_j\sin^2\like{theta}{\symglyph{theta}_j}$, including an infinite right-hand side.\\\addlinespace
\path{DK-4.3-prop} &
 Same displacement setting; put
 $\sym{SquaredDisplacement}(W)=(\sym{Identity}-W)^*(\sym{Identity}-W)$. &
 $\sym{NormN}(\sym{SquaredDisplacement}(\sym{SourceDirectRotation}))
 \le\sym{NormN}(\sym{SquaredDisplacement}(\sym{Wcompetitor}))$:
 the direct rotation minimizes squared full displacement.\\\addlinespace
\path{DK-4.4-prop} &
 Real space; largest principal angle at most $\sym{PiConstant}/3$. &
 The printed assertion that $\sym{SourceDirectRotation}$ minimizes
 every $\sym{NormN}(\sym{Identity}-\sym{Wcompetitor})$ is false.
 A real four-dimensional trace-norm counterexample is given in
 \cref{sec:prop44}; the Q-norm repair is stated there separately.\\\addlinespace[5pt]
\multicolumn{3}{@{}l}{\textbf{Sylvester estimates and spectral cutoffs}}\\*\addlinespace[3pt]
\path{DK-5.1-thm} &
 Banach spaces; bounded $C,X,Y$; $BX-XC=Y$;
 $\dkOperatorNorm{C}\le a$,
 $\dkOperatorNorm{B^{-1}}\le(a+\sym{delta})^{-1}$, $a\ge0$.
 If $B$ is unbounded, it is closed and densely defined with bounded inverse.
 The map norm is compatible with these operator norms. &
 $\sym{delta}\|X\|\le\|Y\|$.
 The coefficient hypotheses may be interchanged.\\\addlinespace
\path{DK-5.2-thm} &
 Self-adjoint $B\ge(\sym{SylvesterGapThreshold}+\sym{delta})\sym{Identity}$,
 $C\le\sym{SylvesterGapThreshold} \sym{Identity}$; bounded $X,Y$;
 $X\DKDom(C)\subset\DKDom(B)$ and $BX-XC=Y$ there. &
 $\sym{delta}\sym{NormN}(X)\le\sym{NormN}(Y)$,
 including unbounded $B,C$.\\\addlinespace
\path{DK-5.1-lem} &
 Orthogonal projections $\sym{StrongApproxProjection}\to \sym{Identity}$ strongly; bounded $K$. &
 $\sym{SingularValue}(K\sym{StrongApproxProjection})\to \sym{SingularValue}(K)$ for each singular-value index $j$.\\\addlinespace[5pt]
\multicolumn{3}{@{}l}{\textbf{Generalized single-angle bounds}}\\*\addlinespace[3pt]
\path{DK-6.1-lem} &
 Each diagonal block of $K$ is dominated by the corresponding block of
 $L$ in every $\sym{NormN}$. &
 The same domination holds for their block-diagonal sums.
 The converse holds if the two $K$ blocks have identical singular values
 (counting multiplicities), as do the two $L$ blocks.\\\addlinespace
\path{DK-6.2-lem} &
 Orthogonal projections $\sym{PinchingProjectionOne},\sym{PinchingProjectionTwo}$; bounded $K$. &
 $\sym{NormN}(\sym{PinchingProjectionOne} K\sym{PinchingProjectionTwo}+
 (\sym{Identity}-\sym{PinchingProjectionOne})K(\sym{Identity}-\sym{PinchingProjectionTwo}))\le\sym{NormN}(K)$.\\\addlinespace
\path{DK-6.1-prop} &
 $\sym{A},\sym{Acurrent}$ have a common domain and bounded difference
 $\sym{H}$;
 both $\sym{SineGap}(\sym{TrialOperator},\sym{Lambda1})$ and
 $\sym{SineGap}(\sym{A1},\sym{Lambda0})$. &
 $\sym{delta}\sym{NormN}(\sin\sym{Theta})\le\sym{NormN}(\sym{H})$.\\\addlinespace
\path{DK-6.1-thm} &
 $\sym{E0}^*\sym{E0}\ge\sym{FrameLowerBoundEpsilon}^2\sym{Identity}$, $\sym{FrameLowerBoundEpsilon}>0$;
 $\sym{SineGap}(\sym{TrialOperator},\sym{Lambda1})$.
 Source and target dimensions may differ. &
 $\sym{delta}\sym{FrameLowerBoundEpsilon}\sym{NormN}(\sin\sym{Theta0})
 \le\sym{NormN}(\sym{R})$.
 The directed sine has the singular values of
 $\sym{Psource}(\sym{Identity}-\sym{Qtarget})$.\\\addlinespace
\path{DK-6.2-thm} &
 Same lower-frame condition; only
 $\sym{SpectralDistanceOperator}(\sym{SpectrumSymbol}(\sym{TrialOperator}),\sym{SpectrumSymbol}(\sym{Lambda1}))
 \ge\sym{delta}$. &
 $\sym{delta}\sym{FrameLowerBoundEpsilon}\|\sin\sym{Theta0}\|_{\mathrm{HS}}
 \le\|\sym{R}\|_{\mathrm{HS}}$.
 If $\sym{R}$ has finite rank, also
 $\sym{delta}\sym{FrameLowerBoundEpsilon}\dkOperatorNorm{\sin\sym{Theta0}}
 \le\dkOperatorNorm{\sym{R}}\sqrt{\sym{RankOperator}\sym{R}}$.\\\addlinespace
\path{DK-6.3-thm} &
 Isometric Rayleigh--Ritz trial data;
 $\sym{OrderedSeparation}_{\sym{delta}}(\sym{TrialOperator},\sym{Lambda1})$.
 The source dimension may be smaller than the target dimension. &
 $\sym{delta}\sym{NormN}(\tan\sym{Theta0})
 \le\sym{NormN}(\sym{R})$.\\\addlinespace
\path{DK-6.3-lem} &
 Bounded $K$; projections $\sym{CompressionProjectionImage},\sym{CompressionProjectionInput}$ of rank at most $n\ge1$;
 $K\sym{CompressionProjectionInput}=\sym{CompressionProjectionImage} K\sym{CompressionProjectionInput}$;
 $\sum_{j\le n}\sym{SingularValue}(K\sym{CompressionProjectionInput})^2$\newline
 ${}>\sum_{j\le n}\sym{SingularValue}(K)^2-\sym{CompressionTolerance}^2$, $\sym{CompressionTolerance}>0$. &
 $\dkOperatorNorm{\sym{CompressionProjectionImage} K(\sym{Identity}-\sym{CompressionProjectionInput})}<\sym{CompressionTolerance}$:
 a nearly maximizing finite-rank compression has small ``leakage''
 (i.e., the norm of this off-diagonal block is below $\sym{CompressionTolerance}$).\\\addlinespace[5pt]
\multicolumn{3}{@{}l}{\textbf{Spectral selection and repulsion}}\\*\addlinespace[3pt]
\path{DK-8.1-thm} &
 Ordered reference gap
 $\sym{OrderedSeparation}_{\sym{delta}}(\sym{ReferenceBlock0},\sym{ReferenceBlock1})$ with endpoint $\sym{alphaEndpoint}$;
 bounded $\sym{H}$, $\sym{H0}=\sym{H1}=0$.
 For the equivalence, $\sym{DimensionOperator}\sym{CrossedSourceDefect}=\sym{DimensionOperator}\sym{CrossedTargetDefect}$. &
 For a chosen reducing target,
 $\sym{Theta}\le\sym{PiConstant}/4$ iff
 $\sym{Lambda0}\le\sym{alphaEndpoint} \sym{Identity}$ and
 $\sym{Lambda1}\ge(\sym{alphaEndpoint}+\sym{delta})\sym{Identity}$; such a reducing target exists.
 With exact frames obtained by applying $\sym{SourceDirectRotation}$ to the reference frames,
 $\sym{ReferenceBlock1}-\sym{alphaEndpoint} \sym{Identity}\le\sym{ComplementCosineBlock}(\sym{Lambda1}-\sym{alphaEndpoint} \sym{Identity})\sym{ComplementCosineBlock}$,
 as quadratic forms, with the lower-block analogue.
 In finite dimensions there are eigenvalue and symmetric-gauge
 repulsion inequalities.\\\addlinespace
\path{DK-8.2-thm} &
 Double-angle sine hypotheses; $\sym{DimensionOperator}\sym{CrossedSourceDefect}=\sym{DimensionOperator}\sym{CrossedTargetDefect}$;
 either bounded $\sym{H}$ with
 $\dkOperatorNorm{\sym{H}}<\sym{delta}/2$, or bounded
 $\sym{R}$ with $\dkOperatorNorm{\sym{R}}<\sym{delta}/2$;
 $\sym{SpectrumSymbol}(\sym{ReferenceBlock0})\subset
 [\sym{betaEndpoint}-\sym{delta}/2,\sym{alphaEndpoint}+\sym{delta}/2]$. &
 The corresponding ambient or directed double-angle sine bound holds
 on the branch with angles strictly below $\sym{PiConstant}/4$.\\
\end{longtable}
\endgroup

\clearpage
\section{Notation correspondence with Davis--Kahan (1970)}
\label{app:dk-notation-correspondence}
This table compares our notation with the corresponding symbols of
Davis and Kahan~\citep[especially \S1]{davis1970rotation}. The correspondences
are by mathematical role, not just by letter. In the original, the reference
block $A_0$ also serves as the residual's comparison operator. Here
$\sym{TrialOperator}$ is a general trial operator, which need not equal that block.
Our additional matching constructions have no direct counterpart in the source.

\begingroup
\small
\setlength{\tabcolsep}{3pt}
\renewcommand{\arraystretch}{1.12}
\begin{longtable}{@{}>{\raggedright\arraybackslash}p{0.43\linewidth}>{\centering\arraybackslash}p{0.12\linewidth}>{\centering\arraybackslash}p{0.16\linewidth}>{\raggedright\arraybackslash}p{0.23\linewidth}@{}}
\caption{Documented correspondence with the source symbols of Davis--Kahan (1970).}
\label{tab:dk-source-notation}\\
\toprule
Mathematical object & DK (1970) & This paper & DK location \\
\midrule
\endfirsthead
\toprule
Mathematical object & DK (1970) & This paper & DK location \\
\midrule
\endhead
\bottomrule
\endfoot
\dkDavisKahanRows
\end{longtable}
\endgroup
